% !TeX spellcheck = ru_RU
% !TEX root = vkr.tex

\section{Эксперимент}
\label{sec:experiment}

\subsection{Цель и методология}

Целью эксперимента является сравнение результатов автоматической сегментации с ручной разметкой,
выполненной биологами лаборатории Гербарий ГБС РАН, и оценка точности морфометрических измерений.

Тестовая выборка включает 9~изображений листьев из гербарных коллекций.
Для каждого изображения биологи предоставили:
\begin{itemize}
    \item эталонную маску, построенную вручную (\texttt{test/mapped/});
    \item субъективную оценку времени, затраченного на разметку одного листа.
\end{itemize}

Те же изображения были обработаны разработанной системой с минимальным числом подсказок
(как правило, один-два положительных клика по телу листа).
Результаты сохранены в директории \texttt{test/result2/}.

\subsection{Метрики качества сегментации}

Точность сегментации оценивается метрикой IoU (Intersection over Union):
\[
    \mathrm{IoU} = \frac{|M_\text{auto} \cap M_\text{ref}|}{|M_\text{auto} \cup M_\text{ref}|}
\]
где $M_\text{auto}$~--- автоматическая маска, $M_\text{ref}$~--- эталонная маска биолога.
Значение IoU$\,\geq 0.85$ принято как удовлетворительный порог для практических морфометрических задач.

\subsection{Результаты измерений}

% TODO: дополнить таблицу данными ручных измерений биологов и временем разметки
В таблице~\ref{tab:measurements} приведены результаты автоматических измерений для девяти листьев.

\begin{table}[H]
\centering
\caption{Результаты автоматических измерений (пиксели)}
\label{tab:measurements}
\begin{tabular}{lrrrrrrrrrr}
\hline
\textbf{Образец} & \textbf{Длина} & \textbf{Ширина} & \textbf{w\,0.125} & \textbf{w\,0.25} & \textbf{w\,0.375} & \textbf{w\,0.5} & \textbf{w\,0.625} & \textbf{w\,0.75} & \textbf{w\,0.875} \\
\hline
121\_170\_0037 & 1132 & 523  & 120 & 240 & 351 & 460 & 507 & 506 & 450 \\
2954\_0111     &  734 & 598  & 453 & 519 & 523 & 527 & 535 & 553 & 583 \\
2954\_0115     & 1143 & 628  & 518 & 600 & 623 & 575 & 443 & 292 & 163 \\
3013\_0032     & 1548 & 930  & 783 & 877 & 924 & 913 & 886 & 874 & 875 \\
3013\_0040     & 2018 & 844  & 622 & 774 & 836 & 768 & 619 & 427 & 199 \\
3275\_0081     &  851 & 539  & 529 & 530 & 519 & 493 & 434 & 356 & 242 \\
3284\_0007     & 1899 & 1132 & 1007 & 1036 & 982 & 945 & 912 & 841 & 663 \\
3726\_0043     &  855 & 610  & 255 & 364 & 429 & 466 & 520 & 570 & 594 \\
l0105          &  444 & 324  & 292 & 323 & 306 & 282 & 274 & 276 & 272 \\
\hline
\end{tabular}
\end{table}

Профиль ширины отражает форму листа: у образца 3013\_0040 прослеживается классическая
эллиптическая форма с максимумом в центре, тогда как у 2954\_0111 ширина практически
не изменяется вдоль оси, что характерно для лингулятной (язычковидной) формы листа.

\subsection{Сравнение с ручной разметкой}

% TODO: заполнить по данным биологов
Для сопоставления автоматических и ручных измерений вычислялось относительное отклонение
по длине и ширине:
\[
    \delta = \frac{|x_\text{auto} - x_\text{ref}|}{x_\text{ref}} \times 100\%
\]

\begin{table}[H]
\centering
\caption{Сравнение автоматических и ручных измерений}
\label{tab:comparison}
\begin{tabular}{lrrrr}
\hline
\textbf{Образец} & \textbf{IoU} & $\delta_\text{длина}$, \% & $\delta_\text{ширина}$, \% & \textbf{Время ручн., мин} \\
\hline
% TODO: данные от биологов
121\_170\_0037 & --- & --- & --- & --- \\
2954\_0111     & --- & --- & --- & --- \\
2954\_0115     & --- & --- & --- & --- \\
3013\_0032     & --- & --- & --- & --- \\
3013\_0040     & --- & --- & --- & --- \\
3275\_0081     & --- & --- & --- & --- \\
3284\_0007     & --- & --- & --- & --- \\
3726\_0043     & --- & --- & --- & --- \\
l0105          & --- & --- & --- & --- \\
\hline
Среднее        & --- & --- & --- & --- \\
\hline
\end{tabular}
\end{table}

\subsection{Апробация}

Прототип системы был продемонстрирован консультанту~--- м.н.с.\ лаборатории Гербарий ГБС РАН
Шкурко~А.~В. По итогам демонстрации качество автоматической сегментации признано
удовлетворительным для дальнейшего развития системы.
Видеозапись демонстрации прилагается к работе.

В ходе апробации были выявлены следующие направления доработки:
\begin{itemize}
    \item поддержка масштабного коэффициента (пересчёт пикселей в мм);
    \item возможность экспорта визуализаций в дополнение к CSV;
    \item расстановка точек вдоль периметра с заданным шагом (требование из технического задания).
\end{itemize}
